% ===== PRÉAMBULE CANONIQUE — à coller VERBATIM en tête de CHAQUE .tex =====
\documentclass[11pt,a4paper]{article}
\usepackage[utf8]{inputenc}\usepackage[T1]{fontenc}\usepackage{lmodern}
\usepackage[french]{babel}
\usepackage{amsmath,amssymb,mathtools}\usepackage{icomma}
\usepackage{siunitx}\sisetup{locale=FR}  % \qty{}{}, \si{}, \ang{}
\usepackage{xcolor}\usepackage{colortbl}
\usepackage{tikz}\usepackage{pgfplots}\pgfplotsset{compat=1.18}
\usepackage[siunitx,europeanresistors]{circuitikz} % circuits (élec.)
\usepackage[most]{tcolorbox}\usepackage{enumitem}\usepackage{array,booktabs}
\usepackage[a4paper,margin=2.2cm]{geometry}
\usetikzlibrary{arrows.meta,calc,angles,quotes,positioning,patterns,%
  backgrounds,shapes.geometric,decorations.pathmorphing,decorations.markings,intersections}
\definecolor{primary}{RGB}{222,49,99}\definecolor{primarydark}{RGB}{150,22,55}
\definecolor{defcol}{RGB}{0,114,178}\definecolor{methcol}{RGB}{52,92,148}
\definecolor{excol}{RGB}{0,135,90}\definecolor{attncol}{RGB}{200,40,55}
\tcbset{boxbase/.style={enhanced,breakable,boxrule=0.9pt,arc=2.5pt,
  left=3mm,right=3mm,top=2.3mm,bottom=2.3mm,fonttitle=\bfseries,coltitle=white}}
% --- boîtes TUTORIEL (noms EXACTS, ne pas renommer) ---
\newtcolorbox{cle}[1][]{boxbase,colback=defcol!6,colframe=defcol,
  title={\textsc{À retenir}\ifx\relax#1\relax\relax\else\textmd{ : \itshape #1}\fi}}
\newtcolorbox{methode}[1][]{boxbase,colback=methcol!7,colframe=methcol,sharp corners,
  title={\textsc{Méthode}\ifx\relax#1\relax\relax\else\textmd{ --- \itshape #1}\fi}}
\newtcolorbox{exemplebox}[1][]{boxbase,colback=excol!5,colframe=excol,
  title={\textsc{Exemple}\ifx\relax#1\relax\relax\else\textmd{ : \itshape #1}\fi}}
\newtcolorbox{piege}[1][]{boxbase,colback=attncol!6,colframe=attncol,
  title={\textsc{Piège}\ifx\relax#1\relax\relax\else\textmd{ : \itshape #1}\fi}}
\newtcolorbox{remarque}[1][]{boxbase,colback=black!4,colframe=black!45,
  title={\textsc{Remarque}\ifx\relax#1\relax\relax\else\textmd{ : \itshape #1}\fi}}
% --- boîtes EXERCICE / CORRIGÉ (noms EXACTS, auto-comptées) ---
\newtcolorbox[auto counter]{exo}[1][]{boxbase,colback=primary!4,colframe=primary,
  title={\textsc{Exercice}~\thetcbcounter\ifx\relax#1\relax\relax\else\textmd{ --- \itshape #1}\fi}}
\newtcolorbox[auto counter]{corrige}[1][]{boxbase,colback=excol!4,colframe=excol,
  title={\textsc{Corrigé}~\thetcbcounter\ifx\relax#1\relax\relax\else\textmd{ --- \itshape #1}\fi}}
\newcommand{\vv}[1]{\vec{#1}}\newcommand{\ds}{\displaystyle}
\newcommand{\R}{\mathbb{R}}\newcommand{\N}{\mathbb{N}}\newcommand{\Z}{\mathbb{Z}}
% =========================================================================
\begin{document}
% THEME_TITLE: Le poids et la chute des corps
\sisetup{per-mode=symbol}

\noindent Le \textbf{poids} est la force par laquelle la Terre attire un corps. C'est lui qui gouverne
la \textbf{chute} : selon qu'il y a ou non de l'air, la chute est \emph{libre} ou \emph{non libre}.

\begin{cle}[poids, chute libre, chute non libre]
\begin{itemize}[leftmargin=1.5em,itemsep=3pt]
  \item \textbf{Poids} : $\ \vv G=m\,\vv g\ $, soit $G=mg$ (en \si{\newton}), vertical, vers le bas,
        avec $g\approx\qty{9,81}{\metre\per\second\squared}$ (souvent arrondi à \qty{10}{\metre\per\second\squared}).
  \item \textbf{Chute libre} (poids seul, sans air) : $\ \vv G=m\vv a \Rightarrow a=\dfrac{mg}{m}=g$.
        L'accélération vaut $g$, \emph{indépendante de la masse}.
  \item \textbf{Chute non libre} (avec résistance de l'air $\vv F_f$, vers le haut) :
        \[ G-F_f=ma \;\Longrightarrow\; a=g-\frac{F_f}{m}<g. \]
\end{itemize}
\end{cle}

\begin{center}
\begin{tikzpicture}[>=Stealth,scale=1]
  \fill[defcol] (0,0) circle (0.28) node[white,font=\small]{$m$};
  \draw[->,line width=1.1pt,methcol] (0,0.28) -- ++(0,1.35) node[above]{$\vv F_f$ (air)};
  \draw[->,line width=1.1pt,attncol] (0,-0.28) -- ++(0,-1.5) node[below]{$\vv G=m\vv g$};
  \draw[->,line width=1.1pt,primary] (-1.7,0.15) -- ++(0,-1.0) node[below]{$\vv a$};
  \node[align=left,font=\small,anchor=west] at (2.2,0.1)
     {$a=g-\dfrac{F_f}{m}$\\[4pt]
      $F_f\nearrow$ avec la vitesse\\[2pt]
      $F_f=G \Rightarrow a=0$ : \emph{vitesse limite}};
\end{tikzpicture}
\end{center}

\begin{methode}[distinguer masse et poids]
\begin{itemize}[leftmargin=1.5em,topsep=2pt,itemsep=1.5pt]
  \item La \textbf{masse} $m$ (en \si{\kilogram}) est une quantité de matière : \emph{invariable} (Terre, Lune, espace).
  \item Le \textbf{poids} $G$ (en \si{\newton}) est une \emph{force} : il dépend du lieu par $g$.
        Sur la Lune $g$ est $\approx 6$ fois plus faible $\Rightarrow$ poids divisé par $\approx 6$, masse inchangée.
\end{itemize}
\end{methode}

\begin{exemplebox}[le parachutiste et la vitesse limite]
Au début du saut, $v$ est faible donc $F_f$ est petite : $a\approx g$, il accélère. Quand $v$ augmente,
$F_f$ grandit, donc $a=g-\tfrac{F_f}{m}$ diminue. Lorsque $F_f=G$, on a $a=0$ : la vitesse devient
\textbf{constante} (vitesse limite). Le mouvement se stabilise non pas parce que la gravité disparaît,
mais parce que l'air l'équilibre exactement ($1^{\text{re}}$ loi).
\end{exemplebox}

\begin{piege}[trois erreurs fréquentes]
\begin{itemize}[leftmargin=1.4em,topsep=2pt,itemsep=1.5pt]
  \item Ne pas confondre \textbf{masse} (\si{\kilogram}) et \textbf{poids} (\si{\newton}) : « peser \qty{70}{\kilogram} » est un abus, $70$ est la masse.
  \item En chute libre, l'accélération \textbf{ne dépend pas} de la masse : dans le vide, plume et marteau tombent ensemble.
  \item En chute non libre, l'accélération est \textbf{toujours} $a<g$ (jamais au-dessus) : $g$ est le maximum, atteint seulement sans air.
\end{itemize}
\end{piege}
\end{document}
